\section{Related work}
\label{sec:related}

Our contribution sits at the intersection of three literatures that have not
previously been connected: psychometric latent-trait models, meteorological
suitability indices, and latent-variable modelling in the environmental sciences.

\paragraph{Item Response Theory.} IRT is the standard framework in psychometrics
and educational measurement for separating the difficulty of a test item from the
ability of a test-taker, using only binary (correct/incorrect) responses
\citep{lord1968,rasch1960,embretson2000}. The two-parameter logistic model writes
the probability of a correct response as
$P(y{=}1) = \sigma\!\left(a_i(\theta_p - b_i)\right)$ for person ability
$\theta_p$, item difficulty $b_i$, and item discrimination $a_i$
\citep{birnbaum1968}. Parameters are classically estimated by marginal maximum
likelihood, integrating over the person-ability distribution with an EM algorithm
and Gauss--Hermite quadrature \citep{bock1981}; \citet{bakerkim2004} survey the
estimation techniques. The classical model treats items as \emph{discrete}, each
carrying its own scalar difficulty $b_i$ estimated from the responses to that
item. Two established lines of work already relax this by letting difficulty
depend on covariates: the linear logistic test model (LLTM), which writes item
difficulty as a \emph{linear} combination of item features \citep{fischer1973},
and the broader family of \emph{explanatory} item response models, which embed
item and person effects in a generalized linear (or nonlinear) mixed model
\citep{deboeck2004}. Our model is a continuous-item member of this explanatory
family. It extends LLTM/explanatory IRT along four axes: (i) the difficulty
covariate is a \emph{continuous physical driver} (e.g.\ wind speed), not a set of
discrete item features; (ii) difficulty is a \emph{smooth} function (spline or
Gaussian process) rather than LLTM's linear form; (iii) it is \emph{anchored to a
physics-derived curve} so its units are physically interpretable; and (iv) it is
applied to a new domain---environmental suitability---in which no two sessions
share an identical item, so the discrete-item view is not merely inconvenient but
inapplicable. Positioned this way, the contribution is an explicit extension of
explanatory IRT to a continuous environmental driver, not a claim that difficulty
has never been modelled as a function.

\paragraph{Suitability and climate indices.} Environmental decision science has a
long tradition of expert-defined indices that map weather to a rating. The Tourism
Climate Index \citep{mieczkowski1985} combines thermal comfort, precipitation,
sunshine, and wind into a single score; the Holiday Climate Index
\citep{scott2016} refines the weighting scheme, and second-generation indices
introduce climatic thresholds \citep{defreitas2008}. Activity-specific suitability
curves used in operational scoring engines are of the same family: expert-weighted
transformations of physical variables. These indices are, by construction,
population- and expert-level: they express how suitable a condition is
\emph{on average}, and none decomposes that rating into an intrinsic-difficulty
term and a skill term. Inverse Suitability takes such an expert curve as the
\emph{prior} for its difficulty function and then identifies the residual
structure that the single curve cannot express.

\paragraph{Latent-variable models in environmental science.} Latent-variable and
random-effects models are widely used in ecology and environmental statistics---
joint species distribution models built on latent factors \citep{warton2015},
occupancy models that separate a latent ``true state'' from an imperfect detection
process \citep{mackenzie2002}, and item-factor-analytic treatments of ordinal
environmental data. The closest in spirit are occupancy models: our skill term
plays a role analogous to a site- or observer-specific detection propensity,
separating a latent property of interest from the process that reveals it. What is
new relative to these is (i) anchoring the latent difficulty to a physics-derived
curve so that it is measured in interpretable units rather than recovered only up
to an arbitrary scale, and (ii) treating the difficulty as a smooth function of a
continuous physical driver rather than a per-site constant.

\paragraph{Forecast verification.} Our promotion metric, the Brier Skill Score,
is standard in forecast verification \citep{brier1950}: it measures the
mean-squared-error improvement of a probabilistic prediction over a reference. We
use it to test whether skill-conditioning earns its added complexity against the
single-curve baseline (Section~\ref{sec:recovery}).
